%\section{Lexical Semantic Change of Mental-Health-Related Concepts}

Lexical semantic change (LSC) refers to shifts in a word's meaning while its grammatical function remains stable, and constitutes a common form of language change \parencite{campbellHistoricalLinguisticsIntroduction2013}. For example, cloud, initially a meteorological term, broadened in usage to refer to internet-based data storage.

\textit{Concept Creep Theory} (CCT) posits that harm-related concepts are particularly prone to LSC. Many harm-related terms, including addiction, bullying, harassment, prejudice, and trauma, appear to have expanded in meaning since at least the 1970s \parencite{haslamConceptCreepPsychologys2016,haslamHarmInflationMaking2020}. Concept creep distinguishes between two forms of LSC that can occur concurrently: vertical creep and horizontal creep. Vertical creep refers to the loosening of definitions to include milder instances, whereas horizontal creep refers to the extension of definitions to encompass qualitatively new phenomena.

Previous studies have examined concept creep across different mental-health-related concepts. \textcite{vylomovaSemanticChangesHarmrelated2021} found that many of the terms they studied, including addiction, bullying, prejudice, harassment, and trauma, displayed both vertical creep and horizontal creep in a corpus of 825,628 psychology article abstracts from 875 journals and, to a lesser extent, in a general corpus combining the Corpus of Contemporary American English (CoCA) and the Corpus of Historical American English (CoHA). \textcite{baesSemanticInflationTrauma2023} found evidence of vertical creep in the word trauma in Vylomova and Haslam's psychology corpus. \textcite{baesSemanticShiftsMental2023} similarly found vertical creep in mental-health-related concepts in the same psychology corpus, including addiction, anger, stress, and worry, as well as in Vylomova and Haslam's general corpus, including addiction, grief, stress, and worry.

Not all findings have aligned with this pattern. \textcite{xiaoHaveConceptsAnxiety2023} found falling vertical creep for anxiety and depression in Vylomova and Haslam's psychology and general corpora, contrary to expectation. This unexpected result may reflect measurement unspecificity, particularly the absence of discourse context and construct distinctions. For example, depression may refer to psychiatric depression, but also to meteorological or economic phenomena. Similarly, anxiety may refer to a nosological category, a disorder construct, or an underlying human experience. Indeed, \textcite{pislRevisitingSemanticSeverity2025} showed that when psychiatric depression and anxiety as human experience were isolated, both displayed upward vertical creep. \textcite{iacobComputationalAnalysisEmergence2026} examined a range of common therapy-speak terms, including toxic, bipolar, psychopath, narcissistic, and triggered, and found mixed results. In an extension of Vylomova and Haslam's psychology corpus, these terms displayed horizontal creep. However, in a Reddit corpus comprising general Reddit comments and comments from psychology-oriented subreddits, results were mixed.

On balance, most mental-health-related concepts appear to have become milder and broader over the last few decades, although evidence for broadening is weaker.

\subsection*{Factors Influencing Concept Creep of Mental-Health-Related Terms}

As noted above, general LSC is ubiquitous and naturally-occuring to some extent \parencite{campbellHistoricalLinguisticsIntroduction2013}. Frequency and polysemy are said to explain most variance in rates of lexical semantic change \parencite{hamiltonDiachronicWordEmbeddings2016}. However, the causes of the disproportionate levels of LSC observed in harm-related concepts remain uncertain.

Proposed explanations include cultural shifts toward greater sensitivity to harm, postmaterialist values, and diminished exposure to adversity \parencite{haslamHarmInflationMaking2020,furediCulturalUnderpinningConcept2016}. In the case of mental-health-related terms, evidence is unclear as to whether the ``creeping'' of mental-health-related concepts reflects changes in official diagnostic criteria, broader sociocultural factors, or a combination of both.

\textcite{fabianoDiagnosticInflationDSM2020} showed that no revision of the Diagnostic and Statistical Manual of Mental Disorders (DSM) from the third edition onward was reliably more inflationary or deflationary overall. However, specific disorders changed significantly. Most notably, ADHD inflated by 18\%, 33\%, and then 17\% in the three revisions following DSM-III. Rates of autism also inflated by 50\% from DSM-III to DSM-III-R, albeit based on a single study, but deflated by 15\% from DSM-IV to DSM-5.

\section{Computational Approaches to Studying Diachronic Lexical Semantic Change}

Since 2018, advances in deep learning have expanded the methodological repertoire for modelling semantic change \parencite{manningHumanLanguageUnderstanding2022}. In particular, they have supported the development of language models that encode words as increasingly sophisticated vector representations. This marks a shift from count-based approaches, where word meaning is inferred from patterns of co-occurrence, toward prediction-based embeddings, where vectors are learned iteratively through a language-modelling objective \parencite{grimmerTextDataNew2022,jm3}.

Building on these methodological developments, \textcite{baesMultidimensionalFrameworkEvaluating2024} proposed the \textit{SIBling} (\textbf{S}entiment, \textbf{I}ntensity, and \textbf{B}readth) framework as a unified and multidimensional approach to LSC. The framework situates CCT within a broader account of LSC and links this conceptual model to computational methods, including both count-based and embedding-based techniques. It distinguishes at least three dimensions along which terms may change over time: vertical drift, horizontal drift, and sentiment.

Vertical drift refers to changes in severity or intensity. A term may become stronger, as in hilarious shifting from cheerful or amusing to extremely funny, or weaker, as in trauma shifting from brain injuries to milder events such as business loss. This dimension maps onto vertical concept creep. Horizontal drift refers to changes in breadth. A term may become narrower, as in doctor shifting from scholar or teacher to primarily denoting a medical professional, or wider, as in cloud shifting from a meteorological term to internet-based data storage. This dimension maps onto horizontal concept creep. Sentiment refers to changes in connotation. A term may acquire a more positive connotation, as in geek shifting from a derogatory term for odd people to someone passionate about a field, or a more negative connotation, as in retarded shifting from a neutral term for intellectual disability to a highly pejorative term. This dimension is roughly equivalent to destigmatisation and stigmatisation.

According to SIBling, these dimensions can be complemented by evaluating shifts in the frequency of target words and in the thematic content of their collocates.

\section{Empirical Gap}

Two gaps motivate the present study. The first concerns the evidentiary basis of existing work. Research on lexical semantic change in mental-health-related concepts has relied heavily on curated---and often the same---corpora. Studies of general language have used \textcite{vylomovaSemanticChangesHarmrelated2021} combined CoCA and CoHA corpus \parencite{baesSemanticShiftsMental2023,baesMultidimensionalFrameworkEvaluating2024,xiaoHaveConceptsAnxiety2023}, while studies of psychology-specific language have repeatedly drawn on \textcite{vylomovaSemanticChangesHarmrelated2021} psychology corpus \parencite{baesMultidimensionalFrameworkEvaluating2024,baesSemanticShiftsMental2023,xiaoHaveConceptsAnxiety2023,iacobComputationalAnalysisEmergence2026}. Other analyses have turned to Reddit data \parencite{kangConvergingRepresentationsAttentionDeficit2025,iacobComputationalAnalysisEmergence2026}.

These corpus choices are consequential because the broader objective is to infer cultural dynamics from lexical change. Curated corpora may reflect shifts in editorial policy, audience orientation, or ideological stance rather than wider changes in public discourse \parencite{pislRevisitingSemanticSeverity2025}. Reddit, similarly, is not a neutral proxy for general discourse, as its user base, like that of many social media platforms, is demographically skewed \parencite{gjurkovicPANDORATalksPersonality2021}. To date, research on lexical semantic change in mental-health-related concepts has not examined general public web discourse.

The second gap concerns the target concepts themselves. Direct evidence on lexical semantic change in ADHD and autism remains sparse, although \textcite{kangConvergingRepresentationsAttentionDeficit2025} showed that ADHD and autism appear to have converged on Reddit. From 2019 onward, their semantic similarity increased, with ADHD and autism becoming more contextually similar than comparison conditions.

The present study addresses these two gaps.

\section{The Current Study}

The present study is guided by three research questions: (i) how has the prevalence of ADHD- and autism-related terminology in general web discourse changed from 2014 to 2026; (ii) how have the semantic profiles of ADHD and autism evolved over this period, relative to non-clinical negative baseline emotion terms; and (iii) how has the thematic framing of ADHD and autism changed over time in public web discourse?

The web data are collected using Common Crawl. We built an economical and reproducible pipeline to extract quality-gated and substantive web content around ADHD and autism, as well as baseline terms.

Methodologically, the study adopts \textcite{baesMultidimensionalFrameworkEvaluating2024} SIBling Framework with two deliberate deviations. First, for measuring intensity and sentiment, the study uses \textcite{mohammadNRCVADLexicon2025} NRC-VAD Lexicon rather than Warriner norms. This dictionary contains more terms, 20,007 compared with 13,915, and has been shown to be substantially more reliable, with reliability of 0.923 compared with 0.823 aggregated across three dimensions \parencite{mohammadNRCVADLexicon2025}. Second, instead of SIBling's top-down, dictionary-based thematic index geared toward measuring pathologisation, the study employs a bottom-up topic-modelling approach to capture emergent framings of the target terms, such as clinical versus identity/neurodiversity narratives, education/workplace accommodation, and social-justice language.

In all, this project contributes by extending concept creep and therapy-speak research to ADHD and autism; shifting evidence from platform-specific \parencite{iacobComputationalAnalysisEmergence2026,kangConvergingRepresentationsAttentionDeficit2025} or curated corpora \parencite{baesMultidimensionalFrameworkEvaluating2024,baesSemanticInflationTrauma2023,baesSemanticShiftsMental2023,xiaoHaveConceptsAnxiety2023} toward general web discourse; implementing and making available an economical, reproducible Common Crawl pipeline designed to extract high-quality general discourse to study diachronic lexical-semantic change in target terms against matched baseline terms; and validating and extending the SIBling framework.

Consistent with the gravity of the literature \parencite{baesSemanticInflationTrauma2023,baesSemanticShiftsMental2023,xiaoHaveConceptsAnxiety2023,iacobComputationalAnalysisEmergence2026,vylomovaSemanticChangesHarmrelated2021}, we hypothesise that ADHD and autism have become more frequent and their collocates milder and broader since 2014. Given the absence of prior findings, the present study adopts a conservative approach and proposes no direction regarding changes in sentiment or thematic content among terms collocating with the target terms.
